%  Преамбула конспекта «Соболевские пространства»
%  Компилировать pdfLaTeX (кодировка utf8, шрифты T2A).

% Latin Modern не имеет T2A-начертаний (нет t2almr.fd), поэтому кириллица идёт
% через cm-super: семейство cmr в кодировке T2A.
\usepackage[T2A]{fontenc}
\usepackage[utf8]{inputenc}
\usepackage[english,russian]{babel}

\usepackage{amsmath,amssymb,amsthm,mathtools}
\usepackage{geometry}
\geometry{a4paper,margin=2.5cm}

\usepackage{tikz}
\usetikzlibrary{arrows.meta,patterns,decorations.pathmorphing,decorations.pathreplacing,shapes.geometric,calc}
\usepackage{pgfplots}
\usepgfplotslibrary{groupplots,fillbetween}
\pgfplotsset{compat=1.18}

\usepackage[most]{tcolorbox}
\tcbuselibrary{breakable}

\usepackage{enumitem}
\usepackage{booktabs}
\usepackage{hyperref}
\hypersetup{
  colorlinks=true,
  linkcolor=blue!55!black,
  citecolor=blue!55!black,
  urlcolor=blue!55!black,
  pdftitle={Соболевские пространства --- конспект}
}

%  Окружения для теорем
\theoremstyle{plain}
% Счётчики независимы: в лекции теоремы, леммы и утверждения нумеруются
% каждое своим рядом (Т1, Т2, ... / Л1, Л2, ...).
\newtheorem{theorem}{Теорема}
\newtheorem{proposition}{Утверждение}
\newtheorem{corollary}{Следствие}
\newtheorem*{corollary*}{Следствие}
\newtheorem{lemma}{Лемма}
\newtheorem{criterion}{Признак}
\newtheorem*{criterion*}{Признак}

\theoremstyle{definition}
\newtheorem{definition}{Определение}
\newtheorem{example}{Пример}

\theoremstyle{remark}
\newtheorem*{remark}{Замечание}
\newtheorem*{reminder}{Напоминание}

\renewcommand{\qedsymbol}{$\blacksquare$}

%  Пометки конспекта
%    \begin{idea} ... \end{idea}  — пояснения, которых не было на лекции
%    \begin{warn} ... \end{warn} — неразборчиво / описка / требует сверки
\newtcolorbox{idea}[1][]{
  breakable,
  enhanced,
  colback=blue!4!white,
  colframe=blue!45!black,
  boxrule=0.4pt,
  left=6pt,right=6pt,top=5pt,bottom=5pt,
  title={Пояснение},
  fonttitle=\bfseries\small,
  coltitle=white,
  attach boxed title to top left={xshift=6pt,yshift=-2pt},
  boxed title style={colback=blue!45!black,boxrule=0pt,arc=1pt},
  #1
}

\newtcolorbox{warn}[1][]{
  breakable,
  enhanced,
  colback=orange!6!white,
  colframe=orange!65!black,
  boxrule=0.4pt,
  left=6pt,right=6pt,top=5pt,bottom=5pt,
  title={Сверить с записью},
  fonttitle=\bfseries\small,
  coltitle=white,
  attach boxed title to top left={xshift=6pt,yshift=-2pt},
  boxed title style={colback=orange!65!black,boxrule=0pt,arc=1pt},
  #1
}

% Отметка начала лекции
\newcommand{\lecture}[2]{%
  \bigskip
  \noindent\hfill\textit{\small #1\ifx\relax#2\relax\else, #2\fi}\\[-0.3em]
  \noindent\rule{\textwidth}{0.4pt}
  \medskip
}

%  Макросы
\newcommand{\R}{\mathbb{R}}
\newcommand{\N}{\mathbb{N}}
\newcommand{\Om}{\Omega}
\newcommand{\clOm}{\overline{\Omega}}
\newcommand{\dOm}{\partial\Omega}
\newcommand{\eps}{\varepsilon}
\newcommand{\dist}{\operatorname{dist}}
\newcommand{\supp}{\operatorname{supp}}
\newcommand{\esssup}{\operatorname*{ess\,sup}}
\newcommand{\Lp}[1]{L^{#1}}
\newcommand{\norm}[1]{\left\lVert #1 \right\rVert}
\newcommand{\abs}[1]{\left\lvert #1 \right\rvert}
\newcommand{\Lpn}[2]{\norm{#1}_{L^{#2}(\Omega)}}
\newcommand{\ball}[2]{B(#1,#2)}
\newcommand{\unif}{\rightrightarrows}
\newcommand{\Lloc}{L^1_{\mathrm{loc}}}
\newcommand{\Wp}[2]{W_{#1}^{#2}}
\newcommand{\Wpz}[2]{\mathring{W}_{#1}^{#2}}
\newcommand{\sgn}{\operatorname{sign}}
\newcommand{\dd}[2]{\frac{\partial #1}{\partial #2}}

% Маркер «++ в кружке» из тетради: доказательство будет дописано позже
% по записям прошлых семестров.
\newcommand{\plusplus}{%
  \tikz[baseline=(pp.base)]{%
    \node[draw, ellipse, inner xsep=3pt, inner ysep=1pt, line width=0.4pt] (pp)
      {\footnotesize $+\,+$};}%
}

%  Вспомогательные функции для рисунков
%    sstep(t)      — гладкая «ступенька» 6t^5-15t^4+10t^3, склеенная с 0 и 1
%    urho(x,rho)   — схематическое усреднение ступеньки u = 1_{[0,1)}
%    bump(x,rho)   — ядро усреднения omega_rho в размерности n = 1
\pgfmathdeclarefunction{sstep}{1}{%
  \pgfmathparse{(#1<0) ? 0 : ((#1>1) ? 1 : (6*(#1)^5-15*(#1)^4+10*(#1)^3))}%
}
\pgfmathdeclarefunction{urho}{2}{%
  \pgfmathparse{sstep((#1+#2)/(2*#2)) - sstep((#1-1+#2)/(2*#2))}%
}
% У края носителя 1-(x/rho)^2 -> 0, и деление 1/(...) переполняет pgfmath
% (предел ~16384) ещё до вызова exp. Поэтому знаменатель зажат снизу через max:
% показатель никогда не превосходит 15 по модулю, exp(-15) ~ 3e-7 — на графике ноль.
\pgfmathdeclarefunction{bump}{2}{%
  \pgfmathparse{(abs(#1)>=0.999*#2) ? 0 : (2.2523/#2)*exp(-1/max(1-((#1)/#2)^2,0.0667))}%
}

%  Быстрая сборка: \showfiguresfalse выключает все TikZ-рисунки.
%  Нужно, когда бесплатный Overleaf не успевает в лимит времени.
\newif\ifshowfigures
\showfigurestrue

\newcommand{\inputfig}[1]{%
  \ifshowfigures
    \input{#1}%
  \else
    \begin{center}\fbox{\small рисунок выключен: \texttt{\detokenize{#1}}}\end{center}%
  \fi
}
